% !TEX root = ../main.tex
\section{Methodology}
\label{sec:methodology}
We extend our prior work~\cite{landscape} to scan additional \gls{LDAP} ports and devise it for improved \gls{LDAP} product detection. We analyze the hosters of \gls{LDAP} and the network segments (for different sizes of hosters) where \gls{LDAP} servers are most likely to be found. Our \gls{TLS} analysis links to the hosters and products we identify. We hence create a new and complementary view of the \gls{LDAP} ecosystem. Note that we do not retrieve any sensitive information from the servers. Below, we provide a description of each improvement we made to our tooling.

\mypara{LDAP scanning and retrieval}
To determine servers with open \gls{LDAP} ports, we use \zmap for all routable prefixes on ports 389, 636, and also 3268 and 3269. The latter two are the ports for \gls{GC}. To the best of our knowledge, \gls{GC} ports have only been investigated in the Rapid 7 blog post of 2016 so far. As \gls{LDAP} over UDP has a restricted use case, and our interest lies in the general use of public \gls{LDAP}, we scan only with TCP. Our campaign runs from 2024-10-29 to 2024-11-01.

We use a custom \gls{LDAP} scanner, written in Go, that supports both \starttls and implicit \gls{TLS} and carries out a so-called \textit{bind} operation to verify that the remote server speaks \gls{LDAP}.
In contrast to our previous work, we now retrieve the Root \gls{DSE} and store the raw data, which allows us to connect it to our new identification stage (see below).
We craft a \gls{LDAP} search request 
%using a filter \textit{(objectClass=*)} and an empty \gls{DN}, 
which limits the scope to the base entry node of the directory tree. This allows us to still retrieve important server attributes without carrying out an authentication process with actual credentials. The attributes we retrieve include vendor name, product version, supported versions of \gls{LDAP}, and a number of vendor-specific attributes that allow us to identify \gls{LDAP} servers that do not identify themselves in other attributes. In particular, it allows us to identify various \gls{LDAP} products by Microsoft ~\cite{ms-ldap-attr}.

\label{subsec:ldap_deploy_host}
We extend Rapid7's \textit{recog} tool~\cite{recog} to analyze the attributes we receive (raw \gls{LDAP} Root DSE entries). We extend it with 13 improved patterns where one (OpenLDAP) has been merged into Rapid'7s source code and 12 others (\gls{AD}-related) are available in our repository. Importantly, this now also allows us to identify more \gls{LDAP} products than prior work and, in many cases, also the operating system. 
To determine product lifecycle status (\eg unsupported versions), we contacted the vendors if no public documentation was available.

\mypara{Networks hosting \gls{LDAP} servers}
We define the size of an \gls{AS} as the number of announced IP addresses from routing information (see \textbf{data sources} below), \ie we aggregate over all its announced prefixes on 2024-10-29. An initial analysis shows that \gls{LDAP} servers are seemingly distributed over networks and prefixes of all kinds and sizes, with no clear pattern discernible in which network segments or prefixes (of some size) of an \gls{AS} they are hosted. To understand whether \gls{LDAP} servers occur more commonly in networks or prefixes of particular sizes, we divide the identified, larger prefixes and networks into slices using the Pareto principle. We take an interval that corresponds to a network of a particular size and divide it into two sub-intervals. The exact position of the dividing line is the power of 2 that is closest to 80\% of the interval size. We keep the 20\% portion as one slice and then iteratively divide the remainder (80\%) again, following the same rule. This results in a list of 20 intervals as shown in Figure~\ref{fig:prefix-len-slices}. The largest slice we consider corresponds to networks that have between $2^{26}$ and $2^{27}$ IP addresses. The smallest slice corresponds to networks with up to $2^8$ IP addresses, \ie a /24 prefix, which is also the smallest routeable prefix length on the Internet. As we show in Section~\ref{sec:results}, this division allows us to derive a clearer characterization of the hosting of \gls{LDAP} servers across networks of various sizes.

\mypara{X.509 certificates and cryptography}
We use a custom tool written in Go to validate the X.509 certificates of \gls{LDAP} servers. This proved necessary as existing tools implemented only a subset of our requirements, in particular fast, offline validation with checks for so-called cross-signing between \glspl{CA} as well as custom inspection of self-signed certificates and intermediate certificates with the \textit{CA} flag set. We use this feature set to investigate possible deployments of private PKIs. We consider the \glspl{CA} in the major root stores: Apple, Java, Microsoft, Mozilla, Chrome, and Google Trust Services.

We perform a security analysis of the public key attributes of the certificates. We check if the public key length complies with the standards set by the CA/Browser forum, \ie we hold the certificates to the same standards that one would expect from Web connections. We check if a public key is flagged as insecure by \textit{badkeys}~\cite{badkeys}, \eg susceptible to known attacks, and if the corresponding private key is known to the online database \textit{pwnedkeys.com}~\cite{pwnedkeys}. 

Our measurements collect \gls{TLS} connection data (implicitly and with \starttls), including cipher suites and protocol versions. We evaluate the use of \gls{TLS} by \gls{LDAP} servers according to RFC 9325, correlating this with their hosters and the products we detect.

\mypara{Data sources}
For the IPs we identify as publicly accessible \gls{LDAP} servers, we determine the corresponding \glspl{AS} with \textit{pyasn} and RIB data from RouteViews. We identify the network type (\gls{ISP}, data center/cloud, etc.) and hoster using the \textit{ip2location} database. Note that one AS may have multiple network types, as this mapping is per prefix.

We used the \gls{CUIDS} from February 2022 to October 2024, with one snapshot per month, to verify the online presence of public-facing \gls{LDAP} servers over time up to the time of our own measurement campaign. Here, we consider the IPv4 addresses that have \gls{LDAP} service running on ports 389, 636, 3268, and 3269.
Our own scans were necessary as Censys does not support \gls{LDAP} \starttls, which restricts our \gls{TLS} analysis, and also does not include raw Root DSE responses, which we need for our product identification.

\mypara{Limitations}
\label{limitations}
Our focus is on publicly accessible \gls{LDAP} servers. Hence, our results do not necessarily generalize to the (possibly much larger) population of \gls{LDAP} servers in internal networks. We conduct our measurement campaign from a single vantage point in Europe.
For geolocation and network type identification, we use the external database \textit{ip2location}, which means we share any bias it may have~\cite{_IPGeolocationCapabilities_2024}.

\gls{LDAP} servers often do not reveal the vendor or product version, limiting our extended \textit{recog} tool when different versions yield identical Root \gls{DSE} responses. This is the case for newer versions of Microsoft's Windows Server (2019 and 2022). We are able to address this to some degree, however, if the server uses TLS 1.3, which is only available from Windows Server 2022 onwards~\cite{alvinashcraft_ProtocolsTLSSSL_2024}. Where another TLS version is used, we cannot distinguish between Windows Server 2019 and 2022. However, we note that both are fairly recent and still supported, unlike quite a few other products we identify in our scans. It is also worth noting that a more intrusive anonymous authentication and attempt to access more sensitive content on the \gls{LDAP} server may have given us further identifying data---however, we refrain from doing this due to ethical concerns. We face similar issues with Kerio Connect and OpenLDAP, neither of which reveals the exact product version.

\mypara{Ethics and Responsible Disclosure}
We followed best practices for Internet measurements as previously documented \cite{ZMapInternetScanner2023}. In particular, we operate an abuse email address. We distribute our scans over a 24-hour period to minimize any impact. No sensitive information was collected. Our scan servers also host a Web server explaining our work, and we use appropriate, meaningful DNS PTR records. We received only a handful of requests to be excluded from our scans.

Following established responsible disclosure practices, we initiated a campaign to notify operators of compromised certificates, of which there are 74 from 11 \glspl{CA}. We used the email addresses listed in the certificate's issuer email attribute for initial contact. We included details such as the serial numbers and a description of our methodology used to identify the weak certificates. As of the time of writing, we have not received any responses. We were unable to find the contact information for six compromised certificates.